\section{Related Work}
\label{sec:related}

\noindent\textbf{Trust-region policy optimization.}
Constraining each update to a neighborhood of the data-generating policy is the
central idea of TRPO~\citep{trpo} and CPO~\citep{cpo}, which enforce an explicit
KL constraint, and of PPO~\citep{ppo}, which approximates it with the cheap
clipped surrogate of Eq.~(\ref{eq:ppo}). Mirror-descent and proximal
formulations~\citep{tomar2022mirror,wang2020truly,wang2019trust} make the
constraint explicit again, and SPO~\citep{spo} replaces the hard clip with a
smooth quadratic penalty whose stationary point lies on the trust-region
boundary. Our work stays within this family but separates a clipping rule into a
proximity criterion and a direction criterion. This separation shows that
existing trust-region masks can improve the proximity criterion while still using
PPO's sampled-ratio direction criterion.

\noindent\textbf{RL for LLM reasoning and off-policy correction.}
PPO-style optimization drives RLHF~\citep{ouyang2022training} and modern reasoning
training~\citep{guo2025deepseekr1,team2025kimi,zhou2025reinforcing,zhou2025variational}. GRPO~\citep{grpo} removes the
critic with group-relative advantages, and recent
follow-ups~\citep{ahmadian2024back,drgrpo,yu2025dapo,cispo} refine the advantage
estimator, the clipping range, or the importance-sampling correction.
DAPO~\citep{yu2025dapo} decouples the upper and lower clip bounds, while
CISPO~\citep{cispo} truncates importance weights instead of masking. These
methods modify the objective or the clipping rule, but they still decide the
update direction from the sampled importance ratio.
In deployed LLM RL systems, an additional source of instability is the
training-inference mismatch. The rollout policy served by an optimized inference
engine can differ from the training policy, which induces an off-policy gap and destabilizes
training~\citep{qi2025defeating,yao2025offpolicy,zheng2025stabilizing}.
Prior work mitigates this gap with truncated or masked importance
sampling~\citep{yao2025offpolicy,zheng2025stabilizing,liu-li-2025,team2025every},
or with engineering-level alignment of numerical and implementation
details~\citep{Team2025EveryAM,he2025nondeterminism}.

\noindent\textbf{Divergence-based trust-region masks.}
Closest to us are methods that replace the sampled ratio deviation $|r-1|$ with a
distributional measure of policy drift. DPPO~\citep{dppo} reinterprets PPO's clip
as a one-sample estimate of total variation and replaces it with an explicit
divergence mask, including the top-$K$ KL realization of Eq.~(\ref{eq:topk-kl})
that we build on.
Complementary work studies non-uniform token-level trust regions that vary the
trust-region constraint across tokens~\citep{mao2026beyond}.
Other divergence-based methods~\citep{r2vpo,yao2026rethinking} use smooth
regularization and reweight updates rather than masking them. All of these
methods inherit PPO's ratio-based direction criterion, which decides whether an
update is outward-pointing from the sign of $\Adv(r-1)$. Our contribution is to
make the direction criterion consistent with the divergence it serves.
We derive the first-order change of the divergence in closed form for softmax
policies and provide truncation-aware estimators of the resulting coefficient.
